\label{mdg:chap:barbero}

La transition d’une hexade à une octade combine trois structures de types
différents : une incidence finie, deux séries de Lambert associées aux
canaux \(q_\pm\) et une identification au paramètre réel de l’action de
Holst. Les trois s’articulent par l’incidence, la fonctionnelle angulaire et
le dictionnaire aréal vers la géométrie de la connexion
d’Ashtekar--Barbero développé dans cette même résidence.

\section{Cardinalités d’incidence 90 et 120}

Soit \(H\) un ensemble de six points et soit \(O\) un ensemble de huit points
contenant une copie distinguée de \(H\). Considérons
\[
 \mathcal F_6=H\times\binom H2,
 \qquad
 \mathcal F_8=O\times\binom H2.
\]
L’inclusion distinguée \(H\hookrightarrow O\) induit
\[
 \iota_{6,8}:\mathcal F_6\hookrightarrow\mathcal F_8,
 \qquad
 (h,\{h_1,h_2\})\longmapsto(h,\{h_1,h_2\}).
\]
Ce morphisme d’incidences est l’antécédent combinatoire des
cardinaux \(90/120\) ; ce n’est ni l’extracteur dodécaphasique ni une série angulaire.

\begin{theorem}[Cardinalités de l’inclusion]
\[
 |\mathcal F_6|=6\binom62=90,
 \qquad
 |\mathcal F_8|=8\binom62=120.
\]
La différence normalisée par vingt-quatre coordonnées est
\[
 \frac{90-120}{24\binom62}=-\frac1{12}.
\]
\end{theorem}
\begin{proof}
Chaque élément de la première coordonnée peut être combiné avec n’importe laquelle des
quinze paires de \(H\). Le principe multiplicatif donne les deux cardinalités ;
la dernière identité résulte de la substitution de \(\binom62=15\).
\end{proof}

\begin{figure}[htbp]
\centering
\begin{tikzpicture}[
  font=\small,
  setnode/.style={circle,draw,minimum size=6mm,inner sep=0pt,
    font=\scriptsize\bfseries},
  box/.style={draw,rounded corners=1.5mm,align=center,
    text width=5.0cm,inner xsep=6pt,inner ysep=5pt},
  arr/.style={-{Stealth[length=2.2mm]},line width=.7pt}]

  % Inclusión geométrica. Las flechas parten del contorno y no atraviesan nodos.
  \node[setnode,draw=hmtblue,fill=hmtlightblue] (h1) at (-2.3,1.25) {1};
  \node[setnode,draw=hmtblue,fill=hmtlightblue] (h2) at (-1.1,1.9) {2};
  \node[setnode,draw=hmtblue,fill=hmtlightblue] (h3) at (.1,1.25) {3};
  \node[setnode,draw=hmtblue,fill=hmtlightblue] (h4) at (.1,-.15) {4};
  \node[setnode,draw=hmtblue,fill=hmtlightblue] (h5) at (-1.1,-.8) {5};
  \node[setnode,draw=hmtblue,fill=hmtlightblue] (h6) at (-2.3,-.15) {6};
  \node[setnode,draw=hmtgreen,fill=hmtlightgreen] (o7) at (-3.2,.55) {7};
  \node[setnode,draw=hmtgreen,fill=hmtlightgreen] (o8) at (1.0,.55) {8};
  \node[draw=hmtblue,dashed,rounded corners=4mm,
    fit=(h1)(h2)(h3)(h4)(h5)(h6),inner sep=4mm,
    label={[hmtblue,font=\bfseries]above:{hexada \(H\), \(|H|=6\)}}] (H) {};
  \node[draw=hmtgreen,line width=.8pt,rounded corners=6mm,
    fit=(H)(o7)(o8),inner sep=5mm,
    label={[hmtgreen,font=\bfseries]below:{octada \(O\), \(|O|=8\), \(H\subset O\)}}] (O) {};

  \node[box,draw=hmtblue,fill=hmtlightblue] (f6) at (5.3,1.35)
    {Incidences définies sur \(H\)\\[1mm]
     \(\mathcal F_6=H\times\binom H2\),\quad
     \(|\mathcal F_6|=6\binom62=90\)};
  \node[box,draw=hmtgreen,fill=hmtlightgreen] (f8) at (5.3,-1.0)
    {Extension des incidences à \(O\)\\[1mm]
     \(\mathcal F_8=O\times\binom H2\),\quad
     \(|\mathcal F_8|=8\binom62=120\)};
  \draw[arr,hmtblue] (O.east) -- ++(.45,0) |- (f6.west);
  \draw[arr,hmtgreen] (O.east) -- ++(.45,0) |- (f8.west);
  \draw[arr,hmtgold] (f6.south) -- node[right,align=left,font=\footnotesize]
    {inclusion \(\iota_{6,8}\)} (f8.north);

  \node[box,draw=hmtgold,fill=hmtlightgold,text width=10.9cm]
    (def) at (1.45,-3.80)
    {Défaut combinatoire orienté\\[-.2mm]
     avec la normalisation déclarée :\\[1mm]
     \(\displaystyle
       \frac{|\mathcal F_6|-|\mathcal F_8|}{24\binom62}
       =\frac{90-120}{24\cdot15}=-\frac1{12}.\)};
  \draw[arr,hmtgold] (f8.south) -- (f8.south |- def.north);

  \node[align=center,text=hmtgray,font=\footnotesize,text width=11.2cm]
    at (1.45,-5.50)
    {Cette identité appartient à l’application d’incidences. Elle n’identifie pas
     l’extracteur dodécaphasique, le semi-groupe de tours ni la fonctionnelle angulaire.};
\end{tikzpicture}

\caption{Inclusion d’un sous-ensemble distingué de six éléments dans un ensemble de
huit éléments et dénombrements d’incidence. Les cardinaux \(90\) et \(120\) ne fusionnent ni l’extracteur
dodécaphasique, ni la fonctionnelle angulaire, ni le semi-groupe des tours.}
\label{mdg:fig:incidencia-seis-ocho-rev9}
\end{figure}

Les indices initiaux de la hiérarchie angulaire sont ainsi associés à deux
espaces d’incidence : 90 configurations sur l’hexade et 120 sur
l’octade ambiante. Le quotient \(-1/12\) est une identité combinatoire de cette
normalisation.

\section{Séries de Lambert associées aux 2 secteurs d’incidence}

Pour \(0<q<1\), nous définissons simultanément
\[
 S_{90}(q)=\frac{q^{90}}{1-q^{270}},
 \qquad
 S_{120}(q)=\frac{q^{120}}{1-q^{360}},
 \qquad
 \Delta S_m=S_m(q_-)-S_m(q_+).
\]
Les développements géométriques
\[
 \Delta S_{90}=\sum_{j\geq0}s_{90+270j},
 \qquad
 \Delta S_{120}=\sum_{j\geq0}s_{120+360j}
\]
situent les deux termes dans le semi-groupe global du
\cref{mdg:chap:torres}.

La fonctionnelle angulaire associée à la transition \(6\to8\) est
\[
 K_{6\to8}=12\Delta S_{90}-\Delta S_{120},
\]
et sa précédence a été fixée par la chaîne fondamentale dodécaphasique :
les douze secteurs et l’unique couture orientée produisent d’abord le poids
\(12:-1\). Le coefficient de \((1-q)^{-1}\), égal à \(1/24\), est ensuite calculé comme sortie exacte ;
l’équation diophantienne qui suit certifie la paire déjà engendrée et ne la
sélectionne pas.
La coordonnée angulaire résultante est
\[
\Gamma_{6\to8}
=\frac{\pi A\Cstar}{180}+K_{6\to8}.
\]
Au point HMT,
\[
\begin{aligned}
 \Delta S_{90}&=2.95169439297457621139847987861\times10^{-4},\\
 \Delta S_{120}&=1.96835112502951720093738706217\times10^{-5},\\
 \frac{\pi A\Cstar}{180}&=0.270485322934140078465586458790\ldots,
\end{aligned}
\]
et, par conséquent,
\[
 \boxed{\Gamma_{6\to8}
 =0.27400767269445927474725526077398041940\ldots.}
\]
Les tableaux calculés avec la coordonnée angulaire conjuguée \(C^*\) arrondie à quinze décimales donnent
\(0.2740076726944593\) en arithmétique à double précision ; la valeur précédente
est l’évaluation multiprécision de la phase régionale complète.

\section{Certificat diophantien ultérieur des poids}

L’identité
\[
 \lim_{q\to1^-}(1-q)\frac{q^m}{1-q^n}=\frac1n
\]
associe à la combinaison \(aS_{90}-bS_{120}\), avant son évaluation
dans les canaux conjugués, le coefficient de \((1-q)^{-1}\)
\[
 \frac a{270}-\frac b{360}.
\]

\begin{theorem}[Poids entiers de coefficient de pôle \texorpdfstring{\(1/24\)}{1/24}]
Les solutions entières positives de
\[
 \frac a{270}-\frac b{360}=\frac1{24}
\]
sont
\[
 (a,b)=(12+3t,1+4t),\qquad t\in\ZZ_{\geq0}.
\]
La paire \((12,1)\) est l’unique solution minimale coordonnée par coordonnée et
l’unique de norme \(\ell^1\) minimale.
\end{theorem}
\begin{proof}
La multiplication par \(1080\) produit \(4a-3b=45\). Modulo trois, on obtient
\(a\equiv0\pmod3\) ; en écrivant \(a=3u\), il en résulte \(4u-b=15\). La
positivité implique \(u=4+t\), \(b=1+4t\), avec \(t\geq0\). Par conséquent,
\(a+b=13+7t\), ce qui démontre la minimalité.
\end{proof}

La primitivité \(\gcd(a,b)=1\) ne sélectionne pas à elle seule la paire minimale, car
il existe d’autres solutions primitives. Le tour fondamental du TPK a déjà
produit le poids \(12:-1\) ; le coefficient de pôle \(1/24\), la positivité et la
minimalité forment son certificat arithmétique ultérieur. Aucun de ces
objets ne se déduit de l’évaluation décimale de \(\Gamma_{6\to8}\).

\newpage
\section{Monotonie de la fonctionnelle par rapport à la coordonnée angulaire conjuguée}

Para \(m\in\{90,120\}\),
\[
 S_m'(q)=\frac{m q^{m-1}(1+2q^{3m})}{(1-q^{3m})^2}>0.
\]
De plus,
\[
 \frac{\dd q_-}{\dd C}=\frac\pi{180}q_-,
 \qquad
 \frac{\dd q_+}{\dd C}=-\frac\pi{180}q_+.
\]

\begin{proposition}[Monotonie stricte]
Pour \(1^\circ\leq C\leq3^\circ\), la fonction
\(C\mapsto\Gamma_{6\to8}(A,C)\) est strictement croissante. Par conséquent,
toute valeur de son image possède au plus un antécédent dans cet intervalle.
\end{proposition}
\begin{proof}
La dérivée de chaque différence \(\Delta S_m\) est positive. Dans l’intervalle
indiqué, on obtient
\(\Delta S'_{120}<5.35\times10^{-4}\), tandis que
\(\pi A/180>0.1273\). Par conséquent,
\(\Gamma'(C)>0.1267\). En \(C=\Cstar\),
\[
 \Gamma'(\Cstar)=0.1328995105618907\ldots.
\]
\end{proof}

L’injectivité locale exclut une seconde coordonnée angulaire proche donnant la
même évaluation de la fonctionnelle. Ce contrôle ne construit ni ne sélectionne
\(C^*\) : la coordonnée angulaire conjuguée a déjà été obtenue dans la branche régionale précédente.
La proposition quantifie uniquement la sensibilité de l’évaluation descendante
une fois \(A\) et \(C^*\) fixés ; elle n’autorise pas à inverser
\(\Gamma_{6\to8}\) ni une valeur physique de \(\gamma_{\mathrm{BI}}\) pour
les redéfinir.
